\documentclass{../../note}

\usepackage{amsthm}
\usepackage{pgfplots}
\pgfplotsset{compat=1.18}
\usepackage{xcolor}
\usepackage{tikz}
\usetikzlibrary{shapes,arrows,positioning,fit,calc,matrix,decorations.pathreplacing}
\usepackage{algorithm}
\usepackage{algpseudocode}
\usepackage{listings}
\usepackage{booktabs}
\usepackage{array}
\usepackage{enumitem}

\title{HUFE 数据库原理模拟试卷 A 参考答案}
\author{isomo}
\setlength{\parindent}{0em}

\begin{document}

\fontsize{12pt}{14.4pt}\selectfont
\maketitle

\section*{一、单项选择题}

1.B\quad 2.A\quad 3.D\quad 4.B\quad 5.B\quad 6.C\quad 7.C\quad 8.C\quad 9.C\quad 10.C\quad 11.B\quad 12.B\quad 13.C\quad 14.B\quad 15.C

\section*{二、填空题}

\begin{enumerate}
\item 模式（概念模式）
\item 数据操作
\item 结构化查询语言
\item 用户定义
\item 读脏数据
\end{enumerate}

\section*{三、判断题}

1.$\surd$\quad 2.$\times$\quad 3.$\surd$\quad 4.$\times$\quad 5.$\surd$\quad 6.$\times$\quad 7.$\surd$\quad 8.$\times$\quad 9.$\surd$\quad 10.$\surd$

\section*{四、简答题}

\begin{enumerate}[itemsep=0.8em]
\item 模式/内模式映像支持物理数据独立性；外模式/模式映像支持逻辑数据独立性。
\item 1NF：属性不可再分；2NF：在 1NF 基础上，非主属性完全函数依赖于每个候选码；3NF：在 2NF 基础上，不存在非主属性对码的传递函数依赖。
\end{enumerate}

\section*{五、综合应用题}

\subsection*{1. SQL 与关系代数题}

\begin{enumerate}[label=(\arabic*), itemsep=0.8em]
\item
\begin{lstlisting}[language=SQL]
CREATE TABLE SC (
    Sno CHAR(10),
    Cno CHAR(10),
    Grade DECIMAL(5,2),
    PRIMARY KEY (Sno, Cno),
    FOREIGN KEY (Sno) REFERENCES Student(Sno),
    FOREIGN KEY (Cno) REFERENCES Course(Cno)
);
\end{lstlisting}

\item
\begin{lstlisting}[language=SQL]
SELECT Sno, Sname
FROM Student
WHERE Sdept = '计算机';
\end{lstlisting}

\item
\begin{lstlisting}[language=SQL]
SELECT S.Sname, SC.Grade
FROM Student S, Course C, SC
WHERE S.Sno = SC.Sno
  AND C.Cno = SC.Cno
  AND C.Cname = '数据库原理';
\end{lstlisting}

\item
\begin{lstlisting}[language=SQL]
SELECT Sname
FROM Student
WHERE Sno NOT IN (SELECT Sno FROM SC);
\end{lstlisting}

\item
\[
\Pi_{Sname}(\sigma_{Cno='C2'}(Student \bowtie SC))
\]
\end{enumerate}

\subsection*{2. 设计与理论题}

\begin{enumerate}[label=(\arabic*), itemsep=0.8em]
\item 联系：教师与班级为 1:N；学员与班级为 M:N。
\item 可转换为：\newline
教师(\underline{工号}, 姓名, 职称)\newline
班级(\underline{班级号}, 班级名, 开课日期, 工号)\quad 外码工号引用教师\newline
学员(\underline{学号}, 姓名, 性别, 联系电话)\newline
选修(\underline{学号}, \underline{班级号})\quad 两个外码分别引用学员和班级。
\end{enumerate}

\end{document}
